\section{Pedestrian Model}\label{sec:model}
% generated by experiments/tables/make_table1.py from data/params/model.yaml; do not edit
\begin{table}[t]
\caption{Model parameters as used by the simulator (top) and the two declared modes (bottom). Values are read from the code and listed in \texttt{data/params/model.yaml}.}
\label{tab:params}
\centering\footnotesize\setlength{\tabcolsep}{3pt}
\begin{tabular}{@{}p{1.75cm}p{4.9cm}p{4.9cm}@{}}
\toprule
Agents & \multicolumn{2}{p{9.9cm}}{$m=80$ kg; $v_0\sim\mathcal N(1.34,0.26)$ m/s, $\ge0.5$; $r\sim\mathcal N(0.25,0.03)$ m, $\ge0.15$; $\tau\sim U(0.4,0.6)$ s} \\
Speed--density & \multicolumn{2}{p{9.9cm}}{$v_{\rm eff}=v_0\,\mathrm{clip}[1-e^{-\gamma(1/\rho-1/\rho_{\max})},0.05,1]$, $\gamma=1.913$, $\rho_{\max}=5.4$ ped/m$^2$; $\rho$ = neighbours within $R=3$ m over $\pi R^2$} \\
SFM & \multicolumn{2}{p{9.9cm}}{$A=2000$ N, $B=0.08$ m, $k=1.2\times10^{5}$ kg/s$^2$, $\kappa=2.4\times10^{5}$ kg/(m\,s)} \\
Walls & \multicolumn{2}{p{9.9cm}}{same law per segment, $A_w=A$, $B_w=B$; variants $B_w$, $B_{w,\rm edge}$ (within 0.5 m of the door edge), rounded edge $r=0.1$ m} \\
TTC & \multicolumn{2}{p{9.9cm}}{$k_{\rm ttc}=1.5$, $\tau_0=3$ s, $\tau_{\max}=8$ s} \\
ORCA & \multicolumn{2}{p{9.9cm}}{$\tau_h=5$ s, $\tau_{\rm ORCA}=0.5$ s, speed cap $2v_0$, gate $w_o=1-\sigma(\rho;4,2)$} \\
Crush & \multicolumn{2}{p{9.9cm}}{$k_{\rm crush}=3.6\times10^{5}$, $\kappa_{\rm crush}=4.8\times10^{5}$, gate $w_c=\sigma(\rho;5.5,3)$} \\
Run & \multicolumn{2}{p{9.9cm}}{$\mathrm{d}t=0.01$ s; speed limit $2v_0$; noise $\sigma$ per component, set per experiment; exit at $|x-g|<0.5$ m; deadlock: all $|v|<0.01$ m/s for $10$ s} \\
\midrule
Mode & published-reproduction & corrected-forward \\
\midrule
Integrator & semi-implicit Euler, then clamp; force clip $\pm10^{6}$ N & clamp, then position update; no clip \\
Spawn & uniform, overlaps allowed & rejection, non-overlapping \\
ORCA & as implemented & RVO2-faithful \\
Crush & present in C4 & off \\
Wall corners & sum over segments & convex-union rule \\
\bottomrule
\end{tabular}
\end{table}

%
\subsection{Forces}
The model is a social force model to which a time-to-collision force and ORCA can be added one at a time (parameters in Table~\ref{tab:params}). Agent $i$ is a disc of radius $r_i$ and mass $m_i$ at $\mathbf{x}_i$ with velocity $\mathbf{v}_i$, desired speed $v_{0,i}$ and relaxation time $\tau_i$. The equation of motion is $m_i\dot{\mathbf{v}}_i=\mathbf{F}_i$ with
\begin{equation}\label{eq:force}
\mathbf{F}_i=\mathbf{F}^{\rm des}_i+\sum_{j}\mathbf{f}^{\rm sfm}_{ij}+\mathbf{F}^{\rm ttc}_i+w_{o,i}\,\mathbf{F}^{\rm orca}_i+\sum_{W}\mathbf{f}^{\rm wall}_{iW},
\end{equation}
where $j$ runs over the other active agents within $R=3$\,m and $W$ over the wall segments. Three configurations switch terms on: SFM keeps the desired, social and wall forces; SFM+TTC adds $\mathbf F^{\rm ttc}_i$; SFM+ORCA adds $w_{o,i}\mathbf F^{\rm orca}_i$.

The drive relaxes the velocity towards a speed $v_{{\rm eff},i}$ along the unit vector $\hat{\mathbf e}_i$ from the agent to its goal $\mathbf g$,
\begin{equation}\label{eq:des}
\mathbf{F}^{\rm des}_i=\frac{m_i}{\tau_i}\bigl(v_{{\rm eff},i}\,\hat{\mathbf e}_i-\mathbf v_i\bigr),\qquad
v_{{\rm eff},i}=v_{0,i}\,\mathrm{clip}\!\left[1-e^{-\gamma\,(1/\rho_i-1/\rho_{\max})},\,0.05,\,1\right].
\end{equation}
Here $v_{{\rm eff},i}$ follows Weidmann's speed--density relation~\cite{weidmann1993}, with $\rho_i=N_i/(\pi R^2)$ and $N_i$ the number of other active agents within $R$. In all \FinTermN\ deadlocks the blocker, the agent nearest the door, has at most \FinTermNbMax\ others within $R$ at detection (median \FinTermNbMedian), so $\rho\le\FinTermRhoMax$\,ped/m$^2$ and $v_{{\rm eff},i}=v_{0,i}$ to within 1\%. The lateral bias $b$, a deterministic tilt of fixed handedness, replaces $\hat{\mathbf e}_i$ in Eq.~\eqref{eq:des} by $(\hat{\mathbf e}_i+b\,\hat{\mathbf e}^{\perp}_i)/|\hat{\mathbf e}_i+b\,\hat{\mathbf e}^{\perp}_i|$, with $\hat{\mathbf e}^{\perp}_i$ the unit normal to $\hat{\mathbf e}_i$.

For a pair at centre distance $d_{ij}$ with $r_{ij}=r_i+r_j$, unit normal $\mathbf n_{ij}$ pointing from $j$ to $i$, tangent $\mathbf t_{ij}$ and $h(x)=\max(0,x)$,
\begin{equation}\label{eq:sfm}
\mathbf f^{\rm sfm}_{ij}=\bigl[A\,e^{(r_{ij}-d_{ij})/B}+k\,h(r_{ij}-d_{ij})\bigr]\mathbf n_{ij}+\kappa\,h(r_{ij}-d_{ij})\,\bigl[(\mathbf v_j-\mathbf v_i)\cdot\mathbf t_{ij}\bigr]\mathbf t_{ij},
\end{equation}
after~\cite{helbing2000nature}. A wall segment acts through its closest point, at distance $d_{iW}$, with the friction taken against the agent's own velocity:
\begin{equation}\label{eq:wall}
\mathbf f^{\rm wall}_{iW}=\bigl[A_w\,e^{(r_i-d_{iW})/B_w}+k\,h(r_i-d_{iW})\bigr]\mathbf n_{iW}-\kappa\,h(r_i-d_{iW})\,(\mathbf v_i\cdot\mathbf t_{iW})\,\mathbf t_{iW}.
\end{equation}
Walls that meet at a corner act as their convex union: at a door edge only the nearest point on the union of jamb and corridor wall acts. The door edges are free, sharp segment ends.

With relative position $\Delta\mathbf x$ and velocity $\Delta\mathbf v$, the contact time $\tau_{ij}$ solves $|\Delta\mathbf x+\Delta\mathbf v\,\tau|=r_{ij}$ and
\begin{equation}\label{eq:ttc}
\mathbf F^{\rm ttc}_i=\sum_j k_{\rm ttc}\,\frac{e^{-\tau_{ij}/\tau_0}}{\tau_{ij}^2}\,\hat{\mathbf n}^{\rm avoid}_{ij},
\end{equation}
where $\hat{\mathbf n}^{\rm avoid}_{ij}$ separates the predicted positions at contact. Pairs count only for $0<\tau_{ij}\le\tau_{\max}$ and $|\Delta\mathbf v|\ge10^{-4}$\,m/s, so $\mathbf F^{\rm ttc}$ vanishes at rest.

\subsection{ORCA as a Force}
ORCA builds one half-plane of permitted velocities per neighbour, with horizon $\tau_h$, and returns the permitted velocity $\mathbf v^{\rm orca}_i$ nearest to a preferred velocity. Two solvers are used: the legacy linear-program solver, which has a leg-normal sign defect (Sect.~\ref{sec:att-lone}), and the RVO2 solver, a line-by-line port of RVO2~\cite{vandenberg2011rvo2}. In force coupling, used throughout Sects.~\ref{sec:physics}--\ref{sec:att-pair}, the preferred velocity is $v_{0,i}\hat{\mathbf e}_i$ (with $v_0$, not $v_{\rm eff}$) and $\mathbf v^{\rm orca}_i$ is turned into a force gated off at high density:
\begin{equation}\label{eq:orca}
\mathbf F^{\rm orca}_i=\frac{m_i}{\tau_{\rm orca}}\bigl(\mathbf v^{\rm orca}_i-\mathbf v_i\bigr),\qquad w_{o,i}=1-s(\rho_i;4,2),
\end{equation}
with $s(x;x_0,c)=1/(1+e^{-c(x-x_0)})$. At rest with no binding constraint $\mathbf v^{\rm orca}_i=v_{0,i}\hat{\mathbf e}_i$, and Eq.~\eqref{eq:orca} gives \emph{ORCA's rest drive}, a constant push $w_{o,i}m_iv_{0,i}\hat{\mathbf e}_i/\tau_{\rm orca}$ towards the goal. ORCA's \emph{drive part} is Eq.~\eqref{eq:orca} with $\mathbf v^{\rm orca}_i$ replaced by $v_{0,i}\hat{\mathbf e}_i$; the remainder, zero when no constraint binds, is its \emph{constraint part}.

\subsection{Integration and Noise}
Because results of social force models depend on the integration scheme~\cite{koster2013}, the update caps the speed before the position update: $\mathbf v'=\Pi(\mathbf v+\mathbf F\Delta t/m)$, $\mathbf v'\leftarrow\Pi(\mathbf v'+\boldsymbol\xi)$, $\mathbf x'=\mathbf x+\mathbf v'\Delta t$, where $\Pi$ caps the speed at $2v_0$. The noise $\boldsymbol\xi$ has independent Gaussian components of standard deviation $\sigma$, either fixed per step or scaled as $\sigma\sqrt{\Delta t/0.01\,{\rm s}}$. In velocity coupling (Sect.~\ref{sec:velorca}) ORCA filters the velocity proposed by the forces and the noise: $\mathbf F$ omits the ORCA term, and after the noise step
\begin{equation}\label{eq:velorca}
\mathbf v'\leftarrow\mathbf v'+w_{o,i}\bigl(\Pi(\mathcal O_i(\mathbf v'))-\mathbf v'\bigr),
\end{equation}
where $\mathcal O_i(\mathbf u)$ is the permitted velocity nearest to $\mathbf u$ under the half-planes of the start of the step.

\subsection{Scenarios and Snapshots}
In the room bottleneck the agents are spawned by rejection sampling, without initial overlaps. A run is deadlocked when the criterion of Table~\ref{tab:params} holds without an exit, and its state is saved as a snapshot, a frozen deadlock. A quench restarts a snapshot under one perturbation for at most 120\,s; all perturbations run on the same snapshots and are compared snapshot by snapshot. Deterministic perturbations stop once all speeds stay below $10^{-3}$\,m/s for 10\,s.

\subsection{Statistics}\label{sec:stats}
All bracketed intervals are 95\%. Proportions carry Wilson intervals~\cite{wilson1927}; paired outcomes are compared by the exact McNemar test~\cite{mcnemar1947} and unpaired counts by Fisher's exact test. Times to free are Kaplan--Meier estimates~\cite{kaplanmeier1958}. The noise level $\sigma_{50}$ at which half the runs complete comes from a logistic fit of completion against $\sigma$, with a bootstrap interval that resamples seeds. Jaccard indices, log rate ratios and paired differences carry percentile bootstrap intervals over snapshots, seeds or runs; CroMa means and log flow ratios carry Student-$t$ intervals.

The overlap of two perturbations is the Jaccard index $J=|S_1\cap S_2|/|S_1\cup S_2|$ of their minority outcome sets $S_1$, $S_2$: the unresolved snapshots or runs at the lone trap and in whole runs, the freed standoffs at the pair. It is compared with its expectation for independent sets of the same sizes by a permutation $p$-value. The rate match of a surrogate is checked on random halves: its parameter is refitted on one half and its rate compared with the resolver's on the other.

A replay displaces every coordinate by a Gaussian of standard deviation $10^{-9}$\,m; five replays per snapshot and perturbation are compared with the unperturbed run on whether and when the blocker is freed. The probes are read jointly as one signature, so no multiplicity correction is applied.
